% titel: Blatt 0, A4: Logarithmus

% Teil I: Werkzeuge, kein Lösungsweg.

Logarithmusgleichungen nach $x$ auflösen sowie Aussagen über
Logarithmusregeln als wahr oder falsch beurteilen.

\begin{werkzeugkasten}[Definition und Regeln]
$\log_a(u)=v \ \Longleftrightarrow\ a^v=u$ für $a>0$,
$a\neq1$, $u>0$; es gilt $\log_a(a)=1$.

\begin{align}
  \log_a(u\,v)&=\log_a u+\log_a v, &
  \log_a\frac{u}{v}&=\log_a u-\log_a v,\\
  \log_a(u^r)&=r\log_a u, &
  \log_a a^r&=r .
\end{align}
\end{werkzeugkasten}

\begin{werkzeugkasten}[Exponentialgleichung lösen]
Gleichungen der Form $a^x=b$ ($a>0$, $a\neq1$, $b>0$) löst man über den
Logarithmus:
\begin{equation}
  x=\log_a b .
\end{equation}
Steht der Exponent komplizierter, etwa $a^{x^2}=b$, so gilt
$x^2=\log_a b$, und evtl. sind beide Vorzeichenlösungen zu beachten.
\end{werkzeugkasten}

\begin{achtung}[Darauf achten]
\begin{itemize}
  \item Logarithmen sind nur für positive Argumente definiert; Definitionsbedingungen beachten.
  \item $\log_a(2a)=2$ gilt \emph{nicht} für alle $a$: es ist
        $\log_a(2a)=\log_a 2+1$, also $=2$ nur, wenn $\log_a 2=1$, d.h.
        $a=2$.
  \item Aussagen mit ``für alle'' sind bereits falsch, wenn \emph{ein}
        Gegenbeispiel existiert.
  \item Beim Auflösen von $10^{x^2}=100$ beide Lösungen $x=\pm\sqrt2$
        berücksichtigen.
\end{itemize}
\end{achtung}
